\section{Решётки и их виды. Частичный порядок в решётке. Булевы алгебры. Примеры.}

\subsection*{Определения и аксиомы решётки}

\paragraph{Определение}
\textbf{Решётка} — это множество $M$ с двумя бинарными операциями $\cap$ (пересечение) и $\cup$ (объединение), удовлетворяющими следующим аксиомам:
\begin{enumerate}
    \item \textbf{Идемпотентность:} $a \cap a = a$, $a \cup a = a$.
    \item \textbf{Коммутативность:} $a \cap b = b \cap a$, $a \cup b = b \cup a$.
    \item \textbf{Ассоциативность:} $(a \cap b) \cap c = a \cap (b \cap c)$, $(a \cup b) \cup c = a \cup (b \cup c)$.
    \item \textbf{Поглощение:} $(a \cap b) \cup a = a$, $(a \cup b) \cap a = a$.
\end{enumerate}

\paragraph{Виды решёток}
\begin{itemize}
    \item \textbf{Дистрибутивная решётка:} если операции распределены друг относительно друга:
    $$a \cup (b \cap c) = (a \cup b) \cap (a \cup c), \quad a \cap (b \cup c) = (a \cap b) \cup (a \cap c).$$
    \item \textbf{Ограниченная решётка:} если существуют \textbf{нуль} $0$ ($\forall a: 0 \cap a = 0$) и \textbf{единица} $1$ ($\forall a: 1 \cup a = 1$).
    \item \textbf{Решётка с дополнением:} ограниченная решётка, в которой для каждого $a$ существует $a'$, такой что $a \cap a' = 0$ и $a \cup a' = 1$.
\end{itemize}

\subsection*{Частичный порядок в решётке}

\paragraph{Определение порядка}
В любой решётке можно ввести отношение частичного порядка $\preceq$:
$$a \preceq b \stackrel{\text{def}}{\iff} a \cap b = a \quad (\text{что эквивалентно } a \cup b = b).$$

\paragraph{Связь с гранями}
Решётку можно определить и через порядок. Если в частично упорядоченном множестве для любых двух элементов существуют:
\begin{itemize}
    \item \textbf{Нижняя грань:} $x = \inf(a, b)$ (наибольшая из нижних границ).
    \item \textbf{Верхняя грань:} $y = \sup(a, b)$ (наименьшая из верхних границ).
\end{itemize}
То такое множество является решёткой, где $a \cap b = \inf(a, b)$ и $a \cup b = \sup(a, b)$.

\subsection*{Свойства дополнения}
В ограниченной \textbf{дистрибутивной} решётке дополнение обладает важными свойствами:
\begin{enumerate}
    \item \textbf{Единственность:} у каждого элемента только одно дополнение.
    \item \textbf{Инволютивность:} $a'' = a$.
    \item \textbf{Законы де Моргана:} $(a \cap b)' = a' \cup b'$ и $(a \cup b)' = a' \cap b'$.
\end{enumerate}

\subsection*{Булевы алгебры}

\paragraph{Определение}
\textbf{Булева алгебра} — это дистрибутивная ограниченная решётка, в которой для каждого элемента существует дополнение.
Она объединяет в себе все свойства операций $\cap, \cup, '$ и констант $0, 1$.

\paragraph{Примеры}
\begin{enumerate}
    \item \textbf{Булеан:} $(2^M; \cap, \cup, \overline{A})$ — множество всех подмножеств с операциями пересечения, объединения и дополнения.
    \item \textbf{Логика:} $(\{0, 1\}; \land, \lor, \neg)$ — классическая двузначная логика. Порядком здесь является импликация.
    \item \textbf{Делители числа:} Пусть $P$ — множество произведений различных простых чисел из набора $T$. Тогда $(P; \text{НОД}, \text{НОК}, \text{ДОП})$ — булева алгебра.
    Здесь $0 = 1$, $1 = \prod_{p \in T} p$, а порядок — отношение «делится на».
\end{enumerate}